\section{Gestion des mises à jour des dépendances}
\label{sec:dependances}
\competence{C4.1.1}{Gérer les mises à jour des dépendances et bibliothèques tierces}{La description du processus de mise à jour des dépendances}
\livrable{La description du processus de mise à jour des dépendances (fréquence, périmètre, type automatique ou manuel).}

\subsection{Veille sur les nouvelles versions}
\todo{Outil retenu : \textbf{Renovate} (natif GitLab), qui surveille en continu les nouvelles versions et ouvre des merge requests de mise à jour ; complété par \texttt{npm outdated}. Décrire la configuration \texttt{renovate.json} (annexe~B).}

\subsection{Fréquence des mises à jour}
\todo{Cadence par type : sécurité (au fil des avis), correctives et mineures (MR Renovate groupées, hebdomadaires), majeures (revue mensuelle planifiée).}

\subsection{Périmètre logiciel concerné}
\todo{Dépendances backend (NestJS, Prisma), frontend (Angular), image Docker de base, outils et actions de la CI. Préciser ce qui est suivi par Renovate et ce qui ne l'est pas.}

\subsection{Type de mise à jour et évaluation d'impact}
\todo{Automatique (patch/mineure : MR Renovate auto-mergée si CI verte) vs manuelle (majeure / breaking change : évaluation d'impact, lecture du changelog, tests ciblés).}

\subsection{Intégration sécurisée}
\todo{Évaluation de l'impact sécurité des dépendances par \textbf{Trivy} + \texttt{npm audit} (équivalents open-source de \textbf{Checkmarx SCA} utilisé en entreprise chez Enedis) ; chaque MR passe par la CI (lint + tests + build) avant fusion sur \texttt{main}. Renvoyer vers l'annexe~B.}
